% SPDX-License-Identifier: LPPL-1.3c
% Copyright (C) 2026 Christophe Jorssen
% Author and Current Maintainer: Christophe Jorssen <christophe.jorssen@gmail.com>
% LPPL maintenance status: maintained
% This file is part of luafemm. See LICENSE and LICENSES.md for its terms.
\section{Harmonic circuits and isolated conductors}
\label{sec:ac-circuits}

A current source fixes a current, not its spatial distribution in a massive
conductor. A voltage source fixes a voltage, not the currents in neighbouring
windings. At nonzero frequency, both questions must be solved with the magnetic
field: skin effect, proximity effect and mutual induction change the answer.
This chapter adds those constraints to the harmonic conventions of
Section~\ref{sec:harmonic}. All amplitudes are \emph{peak} phasors with
$\exp(\mathrm i\omega t)$; displayed powers are cycle averages.

\subsection{Two parallel coils, one isolated bar}
Start with a coarse mesh to choose the layout and check signs. Declare a
circuit at picture level, then give each winding a name. In a planar model,
a complete winding has outgoing and returning legs with opposite signed
turns. Legs with the same |winding| belong to one branch; |type=parallel|
connects the branches to the same voltage. A material choice alone does not
specify an electrical connection.

Here the lower branch has 20 turns and the upper one has 40. The source imposes
$0.1\,\mathrm V$ at $2\,\mathrm{kHz}$ with a $30^\circ$ phase. The central
copper bar is electrically isolated: its total longitudinal current must vanish,
although it supports local eddy currents. The field and branch currents are
computed together. Their division is not a DC resistance ratio.

\luafemmexamplebegin{harmonic-circuits}

% manual example: harmonic-circuits/1

\begin{codeexample}[setup code]
\tikzset{
  ac experiment/.style={femm/problem={frequency=2000,depth=20,
    xmin=-8,xmax=8,ymin=-6,ymax=6,mesh size=.45},
    femm/circuit={name=supply,type=parallel,drive=voltage,voltage=.1,phase=15+15}},
  coil leg/.style={draw=orange!70!black,fill=orange!25},
  /femm/region/coil/.style={material=copper,circuit=supply,winding fill=.6,mesh size=.25},
  isolated bar/.style={draw=teal!70!black,fill=teal!20,
    femm/region={material=copper,conductor=solid,current constraint=zero,mesh size=.15}}}
\end{codeexample}

The following picture uses the preceding definitions.

% manual example: harmonic-circuits/2

\begin{codeexample}[width=0pt]
\begin{tikzpicture}[ac experiment,scale=3.5,
    femm/cache={file=\jobname,mode=auto}]

  \path[femm/region={material=air,mesh size=.15}]
    (-3.1,-.3) rectangle (3.1,.3);

  \filldraw[coil leg,femm/region={coil,winding=lower,turns=-20}]
    (-4.7,-3) rectangle (-3.3,-.5);
  \filldraw[coil leg,femm/region={coil,winding=lower,turns=20}]
    (3.3,-3) rectangle (4.7,-.5);
  \filldraw[coil leg,femm/region={coil,winding=upper,turns=-40}]
    (-4.7,.5) rectangle (-3.3,3);
  \filldraw[coil leg,femm/region={coil,winding=upper,turns=40}]
    (3.3,.5) rectangle (4.7,3);
  \filldraw[isolated bar] (-.8,-2) rectangle (.8,2);
  \draw[dashed,red!70!black,femm/profile={name=gap,samples=101}]
    (-3,0)--(3,0);
  \pic[femm/lines=19,femm/phase=0,femm/arrow spacing=12mm] {femm field};
  \pic[femm/currents={spacing=.45,phase=0}] {femm eddy currents};
  \femmmeasure[name=terminal,quantity=circuit,circuit=supply]
  \femmmeasure[name=lower,quantity=circuit,circuit=supply,winding=lower]
  \femmmeasure[name=upper,quantity=circuit,circuit=supply,winding=upper]
  \femmmeasure[name=losses]
  \node[above,text width=6cm,align=center] at (0,6) {Parallel coils and an isolated copper bar};
  \node[below,font=\small] at (0,-6) {$f=2\,\mathrm{kHz},
    \widehat V=0.1\,e^{\mathrm i30^\circ}\,\mathrm V$};
\end{tikzpicture}
\end{codeexample}

% manual example: harmonic-circuits/3

\begin{codeexample}[width=0pt]
\begin{tikzpicture}
  \begin{axis}[width=8cm,height=6cm,grid=major,no markers,
    xlabel={$x$ (mm)},ylabel={$B_y$ (T)},
    legend style={at={(.5,1.03)},anchor=south}]
    \femmplot[profile=gap,component=By,abscissa=x,representation=real,
      plot options={blue!70!black,thick}]
    \addlegendentry{In phase}
    \femmplot[profile=gap,component=By,abscissa=x,representation=imaginary,
      plot options={red!70!black,dashed,thick}]
    \addlegendentry{Quadrature}
  \end{axis}
\end{tikzpicture}
\end{codeexample}

\luafemmexampleend{harmonic-circuits}

The example uses a reusable |/femm/region/coil| style inside each region's key
list. Declare one |femm/region| capture per path. Finish all geometry and source
declarations before the first field pic, measurement or explicit solve. Then
register profiles and plot their stored samples. Branch measurements use the
same |quantity=circuit| interface as terminal measurements.

Before trusting loss values, refine the conductor mesh relative to the skin
depth $\delta=\sqrt{2/(\omega\mu\sigma)}$ and move the artificial outer boundary.
A small algebraic residual is not evidence of spatial convergence. Constrained
solid conductors add field-response solves; begin with a few branches and a
moderate mesh, then refine the final figure.

\subsection{Sources, branches and region choices}
The existing |/femm/circuit| family is shared by DC and AC. Its declaration is
local to the problem; repeated names are rejected. Keys reset to their defaults
for each circuit declaration. Circuit legs inherit ordinary PGF region styles
within TeX groups. Region keys remain scoped to their path or surrounding scope.

\Needspace{7\baselineskip}
\begin{key}{/femm/circuit/type=\meta{series or parallel} (initially series)}
In harmonic mode, |series| imposes equal complex branch currents and adds their
terminal voltages. |parallel| imposes equal complex voltages and adds currents.
A branch is the set of regions with the same |winding| within this circuit.
Mutual induction is included through the common field solution.
\end{key}
\Needspace{7\baselineskip}
\begin{key}{/femm/circuit/drive=\meta{current or voltage} (initially current)}
Choose the prescribed whole-port quantity. |current| is in amperes and
|voltage| in volts; both default to zero. Set only the driven quantity, leaving
the other at zero. |phase|, initially zero, is its phase in degrees and accepts
PGF mathematical expressions. Thus |drive=voltage,voltage=.1,phase=15+15|
means $\widehat V=0.1\exp(\mathrm i30^\circ)\,\mathrm V$. A zero-voltage
solid ring models a shorted turn. Zero current does not mean zero local current
density in a solid conductor.
\end{key}
\Needspace{7\baselineskip}
\begin{key}{/femm/region/current constraint=\meta{none or zero} (initially none)}
A choice handler for an individual connected solid region. |zero| introduces
an unknown voltage such that $\int_S\widehat J\,dS=0$. It requires positive
frequency, positive material conductivity, and a conforming Delaunay mesh.
It cannot accompany a circuit, an impressed current density or phase, or
|conductor=stranded|. Use separate declarations for disconnected conductors;
a region split into disconnected active pieces by later paint operations is
rejected. Holes inside one connected conductor are allowed.

With |none| and no circuit, the existing solid model retains zero imposed
voltage. It does not infer insulation or a floating electrical connection from
a closed TikZ path. The key is inherited like other region keys and resets
at the end of its TeX group.
\end{key}
\Needspace{7\baselineskip}
\begin{key}{/femm/region/conductor=\meta{auto, solid or stranded} (initially auto)}
|auto| selects stranded interpretation for a named circuit region and solid
interpretation otherwise. This preserves the previous default physical model
while making an explicit |solid| choice effective on circuit legs.

|stranded| homogenises fine wires, suppresses resolved eddy currents and retains
conductivity and |winding fill| for resistance and copper loss. A solid circuit
leg instead resolves its nonuniform current density. Such a leg requires
|turns=1| or |turns=-1|, |winding fill=1|, no |conductor length| override, and
positive frequency and conductivity. Each solid leg must be connected. Planar
coupled problems require positive |depth|. In axisymmetry, constrained solid
rings must stay strictly away from $r=0$.
\end{key}

Planar voltage-driven or parallel branches require closed signed turns
($\sum N=0$) so that their linkage and voltage do not depend on the potential
gauge. An open leg remains usable with prescribed series current, but its
terminal linkage, voltage and impedance are not reported. Conductivity and
conductor lengths must be known for voltage or parallel excitation. Unspecified
conductivity is not silently interpreted as an ideal zero-resistance winding.

\subsection{Measurements and plots}
\Needspace{7\baselineskip}
\begin{key}{/femm/measure/winding=\meta{name} (initially empty)}
For |quantity=circuit| in a harmonic model, restrict the measurement to one
named branch of |circuit|. Empty means the whole port. This key resets for each
measurement; an unknown name or a static branch request raises an error. The
result is saved under |name| and may be used by |\femmvalue| in an annotation
or by |\femmresult| when building a PGFPlots table.
\end{key}
\begin{codeexample}[code only]
\femmmeasure[name=branch,quantity=circuit,circuit=supply,winding=lower]
\femmvalue[component=current_rms]{branch}
\edef\CurrentPhase{\femmresult{branch}{current_phase}}
\end{codeexample}
Available terminal phasors are |current| (A), |voltage| (V), |linkage| (Wb)
and |impedance| ($\Omega$), each with suffix |_real|, |_imaginary|,
|_amplitude| and |_phase|. Phase is in degrees and is undefined for an exactly
zero phasor. Current, voltage and linkage also have |_rms|; impedance does not.
Existing |resistance| is the DC equivalent resistance, |power| is active input
power (W), and |reactive_power| is reactive input power (var). Impedance is
undefined at zero port current. Unavailable components raise an error rather
than fabricate a zero. These conventions also apply to legacy harmonic
series-current winding measurements.

For parallel branches the reported linkage is an equivalent port linkage,
$\lambda_{\rm eq}=R_{\rm eq}\sum_b\lambda_b/R_b$, where
$R_{\rm eq}^{-1}=\sum_bR_b^{-1}$. It satisfies
$\widehat V=R_{\rm eq}\widehat I+\mathrm i\omega\widehat\lambda_{\rm eq}$;
it is not the sum of branch linkages. Select a branch when that is the quantity
needed for a winding diagram.

Field profiles retain |J| and |Jeddy| and their real, imaginary, amplitude,
RMS, phase and snapshot representations. In a constrained solid,
|Jeddy| is the zero-net-current redistribution about the DC transport profile.
It includes the compensating induced voltage, making it invariant under
$A\mapsto A+\hbox{constant}$ in planar geometry. For an isolated conductor it
is the entire current. In an unconstrained solid it retains the previous
zero-imposed-voltage definition. |femm eddy currents| uses that same field.
Regional |joule_loss| integrates the total current; |eddy_loss| integrates the
redistribution. For a uniform solid, their difference is its DC transport
loss. Stranded copper loss uses the fill-adjusted conductivity.

\subsection{Coupled equations and verification}
For each solid leg let $\ell=d$ in a planar problem and $\ell=2\pi r$ in
axisymmetry. Its solved voltage $\widehat V_l$ gives
\[
 \widehat J=\frac{\sigma}{\ell}\widehat V_l
                 -\mathrm i\omega\sigma\widehat A,\qquad
 G_l\widehat V_l-\mathrm i\omega\int_{S_l}\sigma\widehat A\,dS
       =s_l\widehat I_b,\qquad
 G_l=\int_{S_l}\frac{\sigma}{\ell}\,dS.
\]
For an isolated conductor the right side is zero. A branch voltage sums its
signed solid-leg voltages and the stranded contributions
$R\widehat I+\mathrm i\omega\widehat\lambda$. The field mass, load, linkage,
current and loss integrals share quadrature. Response columns of the sparse
complex field operator reduce the terminal equations to a small dense complex
system with scaling, pivoting and a backward-residual check. The recombined
field is checked against the requested true-residual tolerance. Response columns
use a tighter tolerance; the reported GMRES iteration count adds their work.
Very fine meshes may exhaust the per-solve iteration limit; this raises an error
instead of storing an unconverged result.

The native cache stores the real and imaginary potentials, then reconstructs
terminal currents and conductor voltages on restoration. Revision 22 includes
current-constraint semantics and invalidates older numerical snapshots. A
changed terminal drive can reuse a compatible mesh; frozen reads neither solve
nor modify the cache. Source, phase and conductor choices enter the solution
fingerprint.

The regression suite checks the analytic transport-current slab under mesh
refinement, isolated-conductor gauge invariance, voltage/current equivalence,
parallel currents against independently solved unit-current ports, axisymmetric
shorted rings, active and reactive power identities, rejected combinations and
cache reconstruction. The optional xfemm oracle independently meshes and solves
transport-current and isolated-conductor fixtures, comparing complex A/B/H/J,
energy and loss with unchanged 2\% scalar and 4\% derivative gates. The generic format matrix checks identical restored
terminal results under plain LuaTeX, LuaLaTeX and ConTeXt MkIV.

External RLC graphs and prescribed-linkage drives are described in
Section~\ref{sec:electrical-networks}. The supported harmonic |.fem|/|.ans|
subset and its representability checks are described in Section~\ref{sec:interchange}.
General nonlinear harmonic laws, time-domain hysteresis and transient motion
remain outside this harmonic circuit model; the current-fed time formulation
is described in Section~\ref{sec:time}. Captured incremental operating-point workflows retain
their prescribed series-current interface. Their bias is not redefined by a
new AC voltage constraint.
